\thispagestyle{empty}

\begin{center}
    \textbf{АННОТАЦИЯ}
\end{center}

\vspace{10pt}

\noindent\textbf{Студент:} Новичков Дмитрий Евгеньевич, группа~3435.

\noindent\textbf{Руководитель ВКР:} Ведяков Алексей Александрович, доцент ФСУиР, к.т.н., ординарный доцент.

\noindent\textbf{Наименование темы:} Исследование применения обучения с подкреплением для повышения эффективности управления роботами с использованием визуально-языковых моделей.

\noindent\textbf{Кафедра/факультет:} Факультет систем управления и робототехники, Университет ИТМО.

\noindent\textbf{Год защиты:} 2026.

\vspace{8pt}
\noindent\rule{\textwidth}{0.4pt}
\vspace{4pt}

\noindent\textbf{Цель исследования.}
Повышение эффективности управления роботами-манипуляторами при выполнении задач, заданных на естественном языке, за счёт интеграции методов обучения с подкреплением (RL) с визуально-языковыми моделями действий (VLA).

\vspace{6pt}
\noindent\textbf{Решаемые задачи.}
\begin{enumerate}[noitemsep, topsep=2pt]
    \item Провести аналитический обзор современных VLA-моделей и способов их интеграции с RL по работам 2023--2026~гг.\ (RT-2, OpenVLA, $\pi_0$, FLOWER, Gemini Robotics~1.5, EUREKA, RL-VLM-F, STARE-VLA, PLD и др.).
    \item Систематизировать подходы по трём направлениям: RL-дообучение VLA (Residual RL, PLD, STARE-VLA); VLM как генератор вознаграждений (EUREKA, Text2Reward, RL-VLM-F, RoboReward); VLM как высокоуровневый планировщик в связке с RL (SayCan, Embodied-R1).
    \item Разработать методику эксперимента: варьируемые факторы (алгоритм RL, способ генерации награды, архитектура и размер VLM/VLA), фиксируемые параметры, план, метрики и процедуру статистической обработки.
    \item Реализовать программный прототип (RL-дообучение VLA и модуль VLM-генерации вознаграждений) в среде симуляции.
    \item Провести экспериментальное сравнение на стандартных бенчмарках (LIBERO, MetaWorld) с базовыми методами в соответствии с разработанной методикой.
    \item Проанализировать результаты, выявить ограничения и сформулировать рекомендации по выбору архитектуры интеграции VLM+RL.
\end{enumerate}

\vspace{6pt}
\noindent\textbf{Краткая характеристика полученных результатов.}

Реализован прототип с двумя ветками: (1)~основная -- RL-дообучение VLA (residual SAC + стадийное вознаграждение) на LIBERO; (2)~дополнительная -- VLM-кодогенерация награды на MetaWorld.

\textbf{Основной результат} (серии~1--4, около 70 запусков, девять VLA-моделей, 3 seed, 95\% ДИ): OpenVLA-OFT-7B, UniVLA-7B, $\pi_0$-3B, CogACT-8B, SmolVLA-2.2B, $\pi_0$-fast-3B, SpatialVLA-4B, Octo-Small, RT-1-X. После residual SAC + stage RL прирост SR составил от +1,5~п.п.\ (OpenVLA-OFT: 97,1\% $\to$ 98,6\% $\pm$0,3) до +4,1~п.п.\ (RT-1-X: 82,9\% $\to$ 87,0\%); промежуточные: UniVLA 95,0$\to$96,8, $\pi_0$ 94,0$\to$96,1, CogACT 93,0$\to$95,4, SmolVLA 92,5$\to$95,3, $\pi_0$-fast 90,6$\to$93,8, SpatialVLA 88,0$\to$91,5, Octo 86,0$\to$89,9. Установлена сильная обратная корреляция ($\rho = -0{,}97$) между baseline SR и приростом. End-to-end разморозка VLA деградирует ($-$0,6~п.п.). Стадийное вознаграждение даёт +0,8~п.п.\ над разреженным. Результат 98,6\% сопоставим с PLD (99,0\% на LIBERO) при бюджете в 3 раза меньше.

\textbf{Дополнительный результат} (серия~5, MetaWorld): VLM-награда с 3 итерациями рефлексии поднимает SR на push-v3 с 48\% до 85\% ($\rho(R,S)$: 0,12 $\to$ 0,81); средний SR по трём задачам -- 77,2\% (vs 36,7\% у разреженной награды).

Сформулированы 5 рекомендаций по выбору архитектуры интеграции VLM+RL: при наличии предобученной VLA (SR\,$>$\,90\%) -- residual RL с замороженной базой и стадийной наградой; при отсутствии VLA -- VLM-генерация награды с итеративным уточнением; SAC предпочтительнее PPO в residual-постановке. Прототип реализован с открытым исходным кодом (\url{https://github.com/mfclabber/bachelor-thesis}).

\vspace{6pt}
\noindent\textbf{Ключевые слова:} обучение с подкреплением, визуально-языковые модели действий, residual RL, VLM-генерация вознаграждений, LIBERO, MetaWorld, OpenVLA, стадийное вознаграждение, робототехническое манипулирование.

\vspace{6pt}
\noindent\rule{\textwidth}{0.4pt}

\vspace{6pt}
\noindent\textbf{Объём работы:} \pageref{LastPage}~страниц, \ref{fig:last}~рисунков, \ref{tab:last}~таблиц, \printbibcount~источника.

\clearpage
